\FloatBarrier
\section{Experiments}
\label{sec:llm}

\subsection{Experimental setup}
\label{sec:setup}

We run five language models --- DeepSeek-v4-flash, GPT-5.6-luna, Qwen3.8-27B, MiMo-v2.5
and GLM-5.3-flash --- on all four tasks, twice each: once with the interface of
\S\ref{sec:episode} alone and once with two standard tools added (\S\ref{sec:results}).
That is $40$ model--task cells and $800$ episodes. Four scripted strategies are run on the
same instances and reported in Appendix~\ref{app:refs}. All five models are run at one
sampling temperature under one prompt and come from a single price tier, which is enough to
show that a habit is not one model's quirk and not enough to characterise the field; the
prompt is the variable we would vary first, since we fixed one defect in it mid-study that
cut infeasible submissions roughly fivefold.

Every task is run on the same $20$ instances. A seed draws a site, a parameter vector
$\theta \sim P_\Theta$ with its own optimum, rather than a noise realisation, so all
agents face identical problems and every contrast below is paired by site. An agent is a
single \texttt{run(env)} call with access to nothing beyond the interface of
\S\ref{sec:episode}: it is not told $P_\Theta$, the forward model, or that a reference
exists, and regret is computed afterwards by the evaluator from $x^\star_\theta$, which
no agent can read. Within a task all agents share the facility, budget, factor ranges and
cycle length; only the strategy varies.

Feasibility checking is free and unlimited for every agent and a rejected design costs no
budget (\S\ref{sec:checking}). A scripted reference uses that freedom when it is
written, by never proposing a design the facility cannot run; a language model uses it at
run time, through the retry loop of \S\ref{sec:harness}.

\subsection{Domain: agriculture}
\label{sec:domain}

Several domains satisfy the four conditions: microbial evolution, bioprocess scale-up,
materials ageing, dose-finding trials. SlowLab is built so that adding one requires no
environment code (\S\ref{sec:interfaces}). We build the first in controlled-environment
horticulture because the conditions there hold jointly and by nature: a crop occupies its
greenhouse for the better part of a year, plant-to-plant variation is large, running a
treatment means growing and selling a crop, and heat stress destroys fruit that cannot be
recovered. None of this is a dial we set, and a benchmark whose
delay, noise and cost were independent hyperparameters could be tuned until the desired
result appeared. Horticulture also supplies the two things a
domain must bring: a mechanistically modelled response with four decades of published
parameters, and a priced output, so that both $f_\theta$ and the cost of experimenting
come from the literature rather than from us.

Appendix~\ref{app:domain} gives the coefficients with their sources, and
Appendix~\ref{app:defects} names the four that are ours rather than the literature's and
what each is worth. Filling in \S\ref{sec:env}: $f_\theta$ is the reduced state-variable TOMGRO model
\citep{jones1999reduced} and $P_\Theta$ its published between-site parameter spread, so a
seed draws a real site rather than a perturbation we chose. Blocks are climate chambers
and units are independently plumbed hydroponic loops, which is what puts four of the six
factors at chamber level and the other two at loop level. The objective is net profit per
square metre per day, priced with public European coefficients. The same appendix explains
the choice of a mechanistic rather than a Gaussian-process ground truth and the caveats that
follow from it.

\subsection{Tasks}
\label{sec:tasks}

The four are the stages an industrial experimental campaign goes through, and each
isolates a capability the previous one does not require. They differ in the question asked
and in the geometry that makes it hard, never in which of the four conditions applies. The variance components, the cycle length,
the per-unit cost and the commitment structure are identical across all four; only the
question and the geometry change.


\textsc{Sanity} is the control. One factor and sixteen unit-slots make it easy, but its
variance components are those of the other three, since relaxing a condition on one task
would contradict the claim that we do not relax them. Any competent closed loop should
beat random search here, so a method that does not has an interface problem rather than a
reasoning problem.

\textsc{Screen} tests allocation under scarcity. It varies all six management factors the
domain exposes on a budget covering one unreplicated sweep, which forces a choice between
precision and coverage, and two rounds leave no room to iterate out of it. Effect sparsity is not injected: it follows from crop physiology
and the price list, and one factor is rail-bound because at 2024 European electricity
prices supplemental light never repays its energy, so recognising that is part of
screening.

\textsc{Optimise} tests iteration under a control hierarchy. One chamber-level and one
loop-level factor mean a valid design must be a split-plot, so where an agent puts a
treatment matters as much as which treatment it is.

\textsc{Transfer} tests whether an agent modelled the mechanism or the surface. Four
resource factors are optimised at the site's own prices; electricity and gas then both rise
by a factor of $2.2$ and the agent must recommend again \emph{with no further experiments}.
The shock size is the 2021--2022 European energy crisis, over which Eurostat's non-household
gas price reached an all-time high of $0.0867$\,\euro/kWh and non-household electricity
excluding taxes peaked at $0.1987$\,\euro/kWh, both two to three times their pre-crisis
level \citep{eurostat_energy_prices}. It falls only on energy: the tomato price,
transplants, labour and purchased CO$_2$ do not move. Every observation returns a revenue
rate, an energy cost rate and an other-cost rate alongside their combination, to the
scripted and language-model interfaces alike, so an agent that modelled the three
separately can re-solve without running anything and one that modelled only the margin
cannot.

%% Table 2 -- generated by scripts/make_tasks_table.py; do not edit by hand.
\begin{table}[t]
\centering\footnotesize
\setlength{\tabcolsep}{3pt}
\begin{tabular}{@{}llccccccc@{}}
\toprule
\textbf{Task} & \textbf{What the agent must do} & \textbf{Factors} & \textbf{Facility} & \textbf{$R \times$ units} & \textbf{days} & $\mathcal{R}^\star(\varnothing)$ & $\sigma_{\mathrm{tot}}/\sqrt{B}$ & ratio \\
\midrule
\textsc{Sanity} & Close the loop at all & 1 (day temp.) & $4\times4$ & $2 \times 8$ & 420 & 0.0055 & 0.0041 & 1.3 \\
\textsc{Screen} & Find which factors matter & 6 (management) & $8\times3$ & $2 \times 12$ & 420 & 0.0044 & 0.0058 & 0.8 \\
\textsc{Optimise} & Find the optimum, respecting the facility & 2 (mixed level) & $4\times3$ & $3 \times 12$ & 630 & 0.0061 & 0.0028 & 2.2 \\
\textsc{Transfer} & Survive an energy price shock & 4 (resource) & $8\times3$ & $3 \times 24$ & 630 & 0.0041 & 0.0042 & 1.0 \\
\bottomrule
\end{tabular}
\caption{The four tasks. \emph{Facility} is chambers $\times$ loops and \emph{$R\times$units} is rounds times units per round; \emph{mixed level} means one chamber-level and one loop-level factor. \emph{$\mathcal{R}^\star(\varnothing)$} is \eqref{eq:eig} with no data: the most any experiment on the task could remove. \emph{$\sigma_{\mathrm{tot}}/\sqrt{B}$} is the standard error of the mean response a full budget of $B$ units achieves (Appendix~\ref{app:tasks}). \emph{ratio} is the first over the second.}
\label{tab:tasks}
\end{table}


\paragraph{How much each task can reward.} The last two columns of
Table~\ref{tab:tasks} are the ceiling $\mathcal{R}^\star(\varnothing)$ and its ratio to
the standard error of the mean response that a full budget can achieve. Only
\textsc{Optimise} is comfortable, at $2.2$. \textsc{Sanity} sits at $1.3$, and
\textsc{Screen} and \textsc{Transfer} at $0.8$ and $1.0$: on those two the whole of what
instance-tailoring could win is at or below what the noise a full budget leaves behind.
They are the tasks on which every agent in
\S\ref{sec:results}, scripted and otherwise, converts a few per cent of what its design
carried --- the two facts are the same fact --- and, at twenty instances, the tasks whose
five-model ordering does not survive resampling the sites: split-half rank correlation runs
in the same order as this column, from $-0.28$ on \textsc{Screen} to $+0.76$ on
\textsc{Optimise} (Appendix~\ref{app:stability}). What each task supports is therefore
different, and \S\ref{sec:results} reads the two low-ratio tasks only through proportions of
submitted designs and through within-model paired contrasts, never through a leaderboard
position. The cause is not noise we chose. It is that
exposing more factors drives the optimum onto constraints where the remaining factors are
worth little to locate, which Appendix~\ref{app:defects} traces. We report the ratio and
read the results in its light; we stopped using it to admit or exclude tasks, because that
makes the task set a function of the simulator's fine detail.

Four scripted strategies --- random search, a Plackett--Burman screen followed by a
response surface, and Bayesian optimisation at one and two replicates per treatment --- are
released with the environment and run on all four tasks in Appendix~\ref{app:refs}. They
are there so that a reader can see what standard methods do here. They are not the
comparison this paper draws: the ceiling $\mathcal{R}^\star(\varnothing)$ is, because it
does not depend on anyone's implementation choices.
\subsection{Results}
\label{sec:results}

Five models, four tasks, $20$ instances each, run twice: once with the interface of
\S\ref{sec:harness} alone and once with two standard tools added --- a screening design
generator, and a posterior-maximisation routine fitted to whatever observations the agent
has collected. That is forty model--task cells, every one with its full twenty instances.
We report the two halves of \S\ref{sec:eval} in turn, realised regret first and design
and decision quality second, then the two effects that only appear because the benchmark
elicits a recommendation every round and runs a paired tool arm.

\subsubsection{The benchmark is far from saturated}
\label{sec:results:outcome}

\begin{table}[t]
\centering\small
\setlength{\tabcolsep}{5pt}
\begin{tabular}{@{}lccccc@{}}
\toprule
\textbf{Model} & 0-shot & $\bar{\mathcal{R}}\downarrow$ & $\mathcal{R}_{\mathrm{cum}}\downarrow$ & $c\uparrow$ & $\eta\uparrow$ \\
\midrule
\multicolumn{6}{@{}l}{\textsc{Sanity}\quad\footnotesize $\mathcal{R}^\star(\varnothing) = 0.0055$} \\
DeepSeek-v4-flash & 0.0375 & \textbf{0.0279(55)} & 258 & 6\% & 18\% \\
GPT-5.6-luna & 0.0269 & \textbf{0.0131(26)} & 297 & 12\% & 37\% \\
Qwen3.8-27B & 0.0256 & 0.0419(135) & 239 & 8\% & 12\% \\
MiMo-v2.5 & 0.0227 & 0.0265(65) & 252 & 11\% & 19\% \\
GLM-5.3-flash & 0.0182 & 0.0202(72) & 297 & 8\% & 25\% \\
\addlinespace
\multicolumn{6}{@{}l}{\textsc{Screen}\quad\footnotesize $\mathcal{R}^\star(\varnothing) = 0.0044$} \\
DeepSeek-v4-flash & 0.2444 & \textbf{0.1037(195)} & 943 & 5\% & 4\% \\
GPT-5.6-luna & 0.1628 & \textbf{0.0803(93)} & 845 & 6\% & 5\% \\
Qwen3.8-27B & 0.2909 & \textbf{0.0893(121)} & 814 & 6\% & 5\% \\
MiMo-v2.5 & 0.1992 & \textbf{0.0871(86)} & 603 & 7\% & 5\% \\
GLM-5.3-flash & 0.1865 & \textbf{0.0866(84)} & 717 & 7\% & 5\% \\
\addlinespace
\multicolumn{6}{@{}l}{\textsc{Optimise}\quad\footnotesize $\mathcal{R}^\star(\varnothing) = 0.0061$} \\
DeepSeek-v4-flash & 0.0412 & \textbf{0.0365(41)} & 405 & 6\% & 16\% \\
GPT-5.6-luna & 0.0248 & \textbf{0.0110(21)} & 450 & 18\% & 45\% \\
Qwen3.8-27B & 0.0288 & 0.0291(63) & 383 & 8\% & 19\% \\
MiMo-v2.5 & 0.0259 & 0.0267(51) & 508 & 7\% & 21\% \\
GLM-5.3-flash & 0.0245 & \textbf{0.0115(20)} & 437 & 14\% & 46\% \\
\addlinespace
\multicolumn{6}{@{}l}{\textsc{Transfer}\quad\footnotesize $\mathcal{R}^\star(\varnothing) = 0.0041$} \\
DeepSeek-v4-flash & 0.2300 & \textbf{0.0758(100)} & 1896 & 7\% & 5\% \\
GPT-5.6-luna & 0.1635 & \textbf{0.0544(45)} & 1896 & 4\% & 7\% \\
Qwen3.8-27B & 0.2427 & \textbf{0.0711(132)} & 1822 & 9\% & 5\% \\
MiMo-v2.5 & 0.2102 & \textbf{0.0778(133)} & 1976 & 9\% & 5\% \\
GLM-5.3-flash & 0.2022 & \textbf{0.0615(101)} & 2004 & 4\% & 6\% \\
\bottomrule
\end{tabular}
\caption{The four quantities of \S\ref{sec:eval}, for language models without tools, $20$ instances per cell. \emph{0-shot} is the model's own recommendation before it sees any data, and is not one of the four. $\bar{\mathcal{R}}$ \eqref{eq:regret} is realised regret after the full campaign, standard error in the last two places, in bold where it beats that model's own 0-shot; $\mathcal{R}_{\mathrm{cum}}$ \eqref{eq:cumregret} is what the experiments themselves forwent, in the same currency summed over occupied slot-days; $c$ and $\eta$ \eqref{eq:capture} are design and decision efficiency under the cross-validated estimator of \S\ref{sec:estimator}, to be read as ranks. $\eta$ near $100\%$ would mean the recommendation is as good as the design allowed. \textsc{Transfer} is scored at the training prices here so the four tasks are comparable; the post-shock number the task exists for is in \S\ref{sec:results}. No cell reaches $\mathcal{R}^\star(\varnothing)$.}
\label{tab:results}
\end{table}


\paragraph{No arm reaches the ceiling.} $\mathcal{R}^\star(\varnothing)$ is the regret a
perfect reasoner carries having run \emph{no} experiment at all. Not one of the forty arms
reaches it. The bare arms sit between $1.8\times$ and $23.4\times$ above it, widest on the
two tasks with the most factors --- $18$--$23\times$ on \textsc{Screen} and $13$--$19\times$
on \textsc{Transfer} against $2$--$8\times$ on the two smaller ones --- and the scripted
references of Appendix~\ref{app:refs} are closer but also short, at $1.6\times$ to
$22.3\times$. Against the weaker reference of a model's own zero-shot recommendation,
thirty-one of the forty cells end below it and nine end \emph{above}: two or three growing
cycles left the model advising something worse than it would have advised for free. All
nine are on \textsc{Sanity} and \textsc{Optimise}, where the zero-shot recommendation is
already close enough that noisy data can only move it the wrong way. The benchmark is not
saturated, and it is least saturated where the design problem is hardest.
$\mathcal{R}^\star(\varnothing)$ is a ceiling rather than a target: no design we can find
removes more than a third of it in one round (Appendix~\ref{app:oracle}), which is why the
gap to it is read as a scale and the comparisons that follow are made between arms.

\paragraph{The two regrets order agents differently.} Realised regret and cumulative
regret are not substitutes, which is why \S\ref{sec:outcome} reports both: on
\textsc{Sanity} and \textsc{Optimise}, the two tasks whose orderings are stable enough to
read (Appendix~\ref{app:stability}), they rank-correlate at $-0.80$ and $-0.60$ within
task, so the model with the best final answer is among the most expensive campaigns to
have got there. An agent judged on the recommendation alone is free to burn the facility
reaching it.

\paragraph{\textsc{Transfer} costs more after the shock than before it.} The column in
Table~\ref{tab:results} is priced at the training energy costs so the four tasks are
comparable. The quantity the task exists for is the one after: electricity and gas rise by
$2.2$ and the agent must re-solve with no further experiments. Every arm is worse there, by
factors of $1.06$ to $1.98$, and each tooled arm is worse than its bare twin. Two reference
agents isolate what is being tested --- identical in algorithm, budget and observations,
differing only in whether they modelled gross margin or modelled revenue and the two cost
streams separately, they score identically before the shock and differ by $1.47\times$
after it ($0.0605$ against $0.0411$, paired $t = 5.5$ over $40$ sites).

\subsubsection{Designs and decisions fail independently}
\label{sec:results:design}

A campaign can fail because the experiments could not have answered the question, or
because they could and the agent did not read them. \eqref{eq:capture} separates the two,
and both are failing here.

\begin{table}[t]
\centering\small
\setlength{\tabcolsep}{5pt}
\begin{tabular}{@{}lcccccccc@{}}
\toprule
 & \multicolumn{2}{c}{\textsc{Sanity}} & \multicolumn{2}{c}{\textsc{Screen}} & \multicolumn{2}{c}{\textsc{Optimise}} & \multicolumn{2}{c}{\textsc{Transfer}} \\
 & \multicolumn{2}{c}{\footnotesize needs $k\!\ge\!2$} & \multicolumn{2}{c}{\footnotesize needs $k\!\ge\!7$} & \multicolumn{2}{c}{\footnotesize needs $k\!\ge\!3$} & \multicolumn{2}{c}{\footnotesize needs $k\!\ge\!5$} \\[1pt]
\cmidrule(lr){2-3}\cmidrule(lr){4-5}\cmidrule(lr){6-7}\cmidrule(lr){8-9}
\textbf{Model} & \textit{bare} & \textit{+tools} & \textit{bare} & \textit{+tools} & \textit{bare} & \textit{+tools} & \textit{bare} & \textit{+tools} \\
\midrule
\multicolumn{9}{@{}l}{\textbf{(a)} \% of designs that cannot separate the main effects (rank-deficient)} \\[2pt]
DeepSeek-v4-flash & 4 & 10 & 100 & 66 & 57 & 54 & 96 & 66 \\
GPT-5.6-luna & 0 & 0 & 100 & 56 & 9 & 20 & 94 & 71 \\
Qwen3.8-27B & 9 & 6 & 96 & 47 & 51 & 44 & 81 & 54 \\
MiMo-v2.5 & 0 & 0 & 100 & 57 & 90 & 57 & 95 & 65 \\
GLM-5.3-flash & 0 & 6 & 100 & 69 & 64 & 59 & 95 & 66 \\
\midrule
\multicolumn{9}{@{}l}{\textbf{(b)} \% of designs with fewer than $d+1$ distinct treatments --- why (a) happens} \\[2pt]
DeepSeek-v4-flash & 4 & 10 & 100 & 62 & 57 & 54 & 86 & 66 \\
GPT-5.6-luna & 0 & 0 & 100 & 56 & 5 & 12 & 92 & 71 \\
Qwen3.8-27B & 9 & 6 & 89 & 45 & 41 & 44 & 68 & 52 \\
MiMo-v2.5 & 0 & 0 & 100 & 57 & 89 & 54 & 95 & 63 \\
GLM-5.3-flash & 0 & 6 & 100 & 69 & 59 & 56 & 93 & 63 \\
\bottomrule
\end{tabular}
\caption{What the submitted designs look like. Every cell is a percentage of that model's designs on that task, pooled over every round of all $20$ instances; \emph{bare} is the interface alone and \emph{+tools} the same model with the two tools. \textbf{(a)} A design is rank-deficient when its distinct treatments, stacked beside an intercept column, do not reach rank $d+1$, so the $d$ main effects cannot be estimated separately. \textbf{(b)} The immediate cause, and almost the whole of it: $k$ distinct treatments give rank at most $k$, so a design needs $k \ge d+1$ before replication can buy anything. The threshold is the task's, not a fixed number --- \textsc{Sanity} needs two and \textsc{Screen} seven --- and (b) tracks (a) to within a few points on every task. The screening tool repairs much of (a) on the two higher-dimensional tasks and none of it on \textsc{Sanity}, which has one factor and nothing to screen.}
\label{tab:diag}
\end{table}


\paragraph{The designs could not have answered the question.} Unaided, $99\%$ of designs on
\textsc{Screen}, $92\%$ on \textsc{Transfer} and $54\%$ on \textsc{Optimise} are
rank-deficient over the factors the task varies, so the main effects they were commissioned
to separate are not estimable at all (Table~\ref{tab:diag}a, the \emph{bare} columns).
Panel~b gives the cause and nearly all of it: $98\%$, $87\%$ and $50\%$ of those designs
carry fewer distinct treatments than the task has factors plus one. The habit behind the
number is the same everywhere --- the median design uses \emph{two} distinct treatments
replicated across the whole facility, whether the task needs two or seven. These are not
careless designs. All but four of the $1137$ submitted declare a randomisation seed, each
distinct treatment carries a median of three to twelve units depending on the task, and the
allocation respects the control hierarchy; they are designs
that buy precision with the units coverage needed, and under a hard parallelism cap the two
draw on the same physical units. An agent that always buys precision is under-powered
against its own question, which is the opposite of the failure we built the benchmark
expecting. Appendix~\ref{app:oracle} separates this from the alternative reading, that a
low $c$ simply means the task is hard: searching the admissible set with knowledge of
$P_\Theta$ and the forward model finds designs reaching $c$ of $27\%$ to $37\%$, so the
best any agent reaches here is $1.4\times$ to $2.8\times$ short of what one round can buy.
The same search puts the median narrow design well below the median wide one on every task,
which is the form the width diagnosis actually takes: a well-placed pair of treatments can
do well, and a typical one does not.

\paragraph{The decisions then use a small fraction of what those designs allowed.} Decision
efficiency is the half of \eqref{eq:capture} that does not blame the design: it divides the
regret a perfect reasoner would carry \emph{on the design the agent actually submitted} by
the regret the agent carried. The $\eta$ column of Table~\ref{tab:results} runs from $4\%$
to $46\%$ over the twenty bare cells, and splits by task rather than by model --- $22\%$ on
\textsc{Sanity} and $29\%$ on \textsc{Optimise} against $5\%$ on \textsc{Screen} and $6\%$
on \textsc{Transfer}. On the two higher-dimensional tasks the campaigns leave roughly
nineteen twentieths of what they bought on the table.

The two failures are not the same failure. Within task, $c$ and $\eta$ rank-correlate at
$+0.80$ on \textsc{Optimise}, which is the one task whose five-model ordering survives
resampling the twenty sites (Appendix~\ref{app:stability}); the two orderings come apart on
\textsc{Transfer} as well, but there neither is reproducible and we do not read it. A single
score would have reported only their product, and \S\ref{sec:results:findings} shows an
intervention that moves them in opposite directions within a model and paired by instance,
which is the evidence that does not depend on any ordering at all.

\subsubsection{Standard tools improve the design and damage the answer}
\label{sec:results:findings}

\begin{table}[t]
\centering\small
\setlength{\tabcolsep}{4.5pt}
\begin{tabular}{@{}llcccc@{}}
\toprule
\textbf{Task} & \textbf{Model} & $\Delta c$ & $\Delta\eta$ & $\Delta\bar{\mathcal{R}}$ & paired $t$ \\
 & & design & decision & final answer & \\
\midrule
\multicolumn{6}{@{}l}{\textsc{Sanity}} \\
& DeepSeek-v4-flash & -1.1 & -5.3 & +42\% & +1.5 \\
& GPT-5.6-luna & -4.9 & -11.5 & +53\% & +1.0 \\
& Qwen3.8-27B & -1.3 & +3.6 & -22\% & -0.6 \\
& MiMo-v2.5 & -4.8 & +3.0 & -9\% & -0.4 \\
& GLM-5.3-flash & -3.3 & -15.5 & +170\% & +2.2$^\ast$ \\
\addlinespace
\multicolumn{6}{@{}l}{\textsc{Screen}} \\
& DeepSeek-v4-flash & +5.4 & -1.2 & +32\% & +1.8 \\
& GPT-5.6-luna & +5.2 & -1.8 & +45\% & +3.3$^\ast$ \\
& Qwen3.8-27B & +5.9 & -1.1 & +24\% & +1.6 \\
& MiMo-v2.5 & +2.1 & -1.4 & +39\% & +2.9$^\ast$ \\
& GLM-5.3-flash & -0.3 & -1.0 & +25\% & +1.7 \\
\addlinespace
\multicolumn{6}{@{}l}{\textsc{Optimise}} \\
& DeepSeek-v4-flash & +0.7 & +19.8 & -56\% & -4.1$^\ast$ \\
& GPT-5.6-luna & -7.1 & -13.2 & +53\% & +1.3 \\
& Qwen3.8-27B & -0.0 & -5.5 & +40\% & +1.8 \\
& MiMo-v2.5 & -1.7 & +17.8 & -45\% & -2.2$^\ast$ \\
& GLM-5.3-flash & -2.7 & -15.5 & +56\% & +1.1 \\
\addlinespace
\multicolumn{6}{@{}l}{\textsc{Transfer}} \\
& DeepSeek-v4-flash & +1.8 & -1.1 & +26\% & +1.2 \\
& GPT-5.6-luna & +6.5 & -3.0 & +59\% & +4.5$^\ast$ \\
& Qwen3.8-27B & -0.2 & -0.7 & +16\% & +1.0 \\
& MiMo-v2.5 & -0.5 & -0.6 & +16\% & +0.9 \\
& GLM-5.3-flash & +4.6 & -0.5 & +4\% & +0.2 \\
\addlinespace
\bottomrule
\end{tabular}
\caption{What the two tools do, paired by instance, on three of the four quantities of Table~\ref{tab:results}. $\Delta c$ is the change in design efficiency in percentage points, $\Delta\eta$ the change in decision efficiency also in percentage points, and $\Delta\bar{\mathcal{R}}$ the relative change in the regret of the final recommendation. A positive $\Delta c$ with a positive $\Delta\bar{\mathcal{R}}$ is a better design with a worse answer, which is what most cells show. $^\ast$ marks $p<0.05$ on a paired $t$-test over instances: six cells, four of them a worse answer and two a better one, both on \textsc{Optimise}.}
\label{tab:tools}
\end{table}


This is
the sharpest evidence we have that the two halves of \eqref{eq:capture} measure different
things, because it is one intervention applied within a model and paired by instance rather
than a correlation across agents --- two numbers that moved together would be one number.
The screening tool works: rank-deficiency falls from $99\%$ to $59\%$ on \textsc{Screen}
and $92\%$ to $64\%$ on \textsc{Transfer}, and $c$ rises in seven of those ten cells by up
to $6.5$ points, the other three moving by $-0.5$ to $-0.2$, inside the estimator's noise.
In the same ten cells $\eta$ falls and realised regret rises, by $4\%$ to $59\%$; over all
four tasks the tools make $16$ of $20$ pairs worse and four better, six significant at
paired $p<0.05$ --- four worse and two better, both on \textsc{Optimise}
(Table~\ref{tab:tools}).

The mechanism is that the inference tool inherits whatever the design can support. It is
the same routine our Bayesian-optimisation reference uses for its own recommendation, and
that reference carries $0.0172$, $0.0586$, $0.0120$ and $0.0335$ on the four tasks against
$0.0711$, $0.1257$, $0.0219$ and $0.0901$ for the identical routine fitted to a language
model's design. Nothing about the routine changed, only what it was fitted to, and what it
was fitted to is thinner: on the three multi-factor tasks the reference's campaigns hand it
$1.8$ to $3.7$ times as many distinct points. Full rank is not sufficiency: the median
tooled design on \textsc{Screen} is thirteen points over six factors, barely one level a
factor, and its main effects are estimable long before its surface is. Appendix~\ref{app:refs} rules out the obvious alternative, that our tool is simply
misconfigured. The models then absorb roughly a third of the tool's disadvantage: regressing
the tools' effect on how much worse the tool's own proposal is than the model's gives a
slope of $0.29 \pm 0.09$ at $r = 0.62$, $p = 0.004$, with the sign predicted in eighteen of
twenty cells. A third is an average over cells, and it hides how literal the deference becomes round by
round: in $384$ of the $693$ turns where both are visible --- $55\%$, and between $53\%$ and
$62\%$ on every task --- the recommendation the model states is the tool's proposed point to
six digits. The models are neither credulous nor immune, and none of them says the tool is
wrong --- which is the regime's doing rather than the tool's. Detecting a bad tool is
ordinarily cheap: run its recommendation, see the number, stop using it. Here that costs a
growing cycle and the agent has at most three, so what the benchmark measures is not whether
a model can use a tool but whether it can price one it cannot afford to test.

\subsubsection{The first cycle of data often makes the recommendation worse}
\label{sec:results:trajectory}

Eliciting a recommendation after every round rather than once at the end exposes something the endpoint
hides. In $44\%$ of episodes the recommendation after the first round is worse than the one
the same model made with no data at all, and in $24$ of the $40$ cells the mean is worse.
The split is by task: $68\%$ of \textsc{Optimise} episodes and $54\%$ of \textsc{Sanity},
against $29\%$ and $26\%$ on \textsc{Screen} and \textsc{Transfer}, where the zero-shot
recommendation starts far enough from the optimum that one round of almost any data helps.
Slow feedback costs an agent in proportion to how many cycles it needs, and here the first
cycle is often spent undoing itself.
